\section{Reproducing the results}
\label{app:reproducing}

Everything in this report comes from the repository. The README describes the setup; this appendix lists what produces each result.

\paragraph{Programs.} The first CMake configure fetches and builds a pinned ns-3 under \file{third_party/}. All programs land in \file{build/bin/}.
\begin{itemize}
  \item \file{spraypoint-verifier}: validates Spraypoint against the paper (\appref{app:validation}).
  \item \file{bypass-model}: the flow-level checker (\secref{sec:checker}).
  \item \file{tcp-experiment} and \file{tcp-bypass}: the ns-3 simulation without and with the bypass.
\end{itemize}

\paragraph{Experiments.} Each experiment is a script in \file{experiments/}, which runs its programs and writes one folder per run under \file{experiments/results/<machine>/<batch>/<script>/}, named after the code commit that produced it. \file{REPRODUCE.sh} runs everything, and \file{REMOTE.sh} runs it on a server over SSH. Rerunning a script skips finished runs.

\paragraph{Figures and tables.} \file{experiments/analysis/report_figures.py} reads those batches and writes every figure in \file{report/figures/} and every generated table in \file{report/tables/}. The notebooks in \file{experiments/analysis/} show the same analyses interactively: \file{01-baseline} and \file{02-bypass-comparison} for any sweep (set \texttt{SWEEP}), and \file{03-model-check} for \file{05} and \file{06} together.

\begin{verbatim}
cmake -B build && cmake --build build --parallel
./REMOTE.sh fetch                     # copy the server's results
experiments/analysis/.venv/bin/python experiments/analysis/report_figures.py
cd report && latexmk             # writes build/report.pdf
\end{verbatim}
